%! Parametrization of Implicitly Defined Functions — Unit 4, Lesson 4.7 — Group Activity (Answer Key)
\documentclass[10pt]{article}
\usepackage{precalculus-article}
\usepackage{precalculus-key}   % pulls in -boxes; do NOT also load -boxes
\newcommand{\TallMath}[1]{$\displaystyle #1\rule[-1.4em]{0pt}{3.2em}$}

\begin{document}

\pageheader{Unit 4, Lesson 4.7}{Group Activity --- Answer Key}
\namepartnerperiod

\vspace{0.1in}

\begin{scenariobox}[Mission Control --- Parametrizing Conic Paths]{sky}
\small
Mission Control is tracking three objects. Each follows a conic section path described by
an implicit equation. Your job: write a parametrization for each object, then use parametric
tools to extract key position data.
\end{scenariobox}

\vspace{0.1in}

%----------------------------------------------------------------------
% Object 1 — Space Station (ellipse)
%----------------------------------------------------------------------
\begin{tcolorbox}[colback=white, colframe=black!40,
    title=\textbf{Object 1 --- Space Station (Ellipse) --- Key},
    fonttitle=\bfseries, arc=2mm, left=3mm, right=3mm, top=2mm, bottom=2mm]
\small
The station's orbital path satisfies $\dfrac{(x-1)^2}{64}+\dfrac{(y+2)^2}{36}=1$.

\textbf{(a)} $h = $ \ans{$1$} \quad $k = $ \ans{$-2$} \quad $a = $ \ans{$8$} \quad $b = $ \ans{$6$}

\vspace{0.08in}
\textbf{(b)} Parametrization:

$x(t) = $ \ans{$1 + 8\cos t$} \qquad $y(t) = $ \ans{$-2 + 6\sin t$}

\vspace{0.08in}
\textbf{(c)} Four extreme positions:

\renewcommand{\arraystretch}{1.6}
\begin{tabular}{lll}
Rightmost ($x$ max): $t = $ \ans{$0$} & Point: \ans{$(9,\,-2)$} \\
Leftmost ($x$ min): $t = $ \ans{$\pi$} & Point: \ans{$(-7,\,-2)$} \\
Highest ($y$ max): $t = $ \ans{$\pi/2$} & Point: \ans{$(1,\,4)$} \\
Lowest ($y$ min): $t = $ \ans{$3\pi/2$} & Point: \ans{$(1,\,-8)$} \\
\end{tabular}

\vspace{0.08in}
\textbf{(d)} $y$-axis crossing: set $x(t)=0 \Rightarrow 1+8\cos t=0 \Rightarrow \cos t = $ \ans{$-\tfrac{1}{8}$}

$t = $ \ans{$\arccos(-\tfrac{1}{8}) \approx 1.696$ rad (and $2\pi - 1.696 \approx 4.587$ rad)}

$y$-axis crossing coordinates: \ans{$\bigl(0,\,-2+6\sin(\arccos(-1/8))\bigr) \approx (0,\,3.97)$ and $(0,\,-7.97)$}

\begin{teachernote}
  Students often stop at $\cos t = -1/8$ and forget to find the $y$-coordinate.
  Prompt: ``You found the time $t$ --- now plug it back into $y(t)$ to get the position.''
  For the exact form: $\sin(\arccos(-1/8)) = \pm\sqrt{1-1/64} = \pm\sqrt{63}/8 = \pm 3\sqrt{7}/8$.
\end{teachernote}
\end{tcolorbox}

\vspace{0.1in}

%----------------------------------------------------------------------
% Object 2 — Comet (parabola)
%----------------------------------------------------------------------
\begin{tcolorbox}[colback=white, colframe=black!40,
    title=\textbf{Object 2 --- Comet (Parabola) --- Key},
    fonttitle=\bfseries, arc=2mm, left=3mm, right=3mm, top=2mm, bottom=2mm]
\small

\textbf{(a)} Vertex: \ans{$(1,\,-3)$} \quad Orientation: \ans{vertical (opens upward)}

\vspace{0.08in}
\textbf{(b)} $x(t) = 1+t$, \quad $y(t) = $ \ans{$2t^2 - 3$}

\vspace{0.08in}
\textbf{(c)} $x$-axis: set $y(t) = 2t^2 - 3 = 0$:

$t = $ \ans{$\pm\sqrt{3/2} = \pm\tfrac{\sqrt{6}}{2} \approx \pm 1.225$}

$x = 1 + t$: \quad \ans{$x = 1 + \tfrac{\sqrt{6}}{2} \approx 2.225$ \quad and \quad $x = 1 - \tfrac{\sqrt{6}}{2} \approx -0.225$}

Coordinates: \ans{$\bigl(1+\tfrac{\sqrt{6}}{2},\,0\bigr)$ and $\bigl(1-\tfrac{\sqrt{6}}{2},\,0\bigr)$}
\end{tcolorbox}

\vspace{0.1in}

%----------------------------------------------------------------------
% Object 3 — Satellite (circle)
%----------------------------------------------------------------------
\begin{tcolorbox}[colback=white, colframe=black!40,
    title=\textbf{Object 3 --- Satellite (Circle) --- Key},
    fonttitle=\bfseries, arc=2mm, left=3mm, right=3mm, top=2mm, bottom=2mm]
\small

\textbf{(a)} Center: \ans{$(0,0)$} \quad Radius: \ans{$10$}

\vspace{0.08in}
\textbf{(b)} $x(t) = $ \ans{$10\cos t$} \qquad $y(t) = $ \ans{$10\sin t$}

Domain: \ans{$0 \le t \le 2\pi$}

\vspace{0.08in}
\textbf{(c)} At $t = \pi/4$:

$x(\pi/4) = 10\cos(\pi/4) = $ \ans{$10 \cdot \dfrac{\sqrt{2}}{2} = 5\sqrt{2} \approx 7.07$}

$y(\pi/4) = 10\sin(\pi/4) = $ \ans{$10 \cdot \dfrac{\sqrt{2}}{2} = 5\sqrt{2} \approx 7.07$}

Coordinates: \ans{$(5\sqrt{2},\,5\sqrt{2})$}

\begin{teachernote}
  A common error: students write $r = 100$ (confusing $r^2$ with $r$). Remind them:
  $x^2 + y^2 = r^2$, so $r^2 = 100$ gives $r = 10$.
\end{teachernote}
\end{tcolorbox}

\end{document}
